%% The first command in your LaTeX source must be the \documentclass command.
%%
%% Options:
%% twocolumn : Two column layout. Do not use twocolumn for papers submitted to CEUR-WS!
%% hf: enable header and footer.
\documentclass[
% twocolumn,
% hf,
]{ceurart}

%%
%% One can fix some overfulls
\sloppy

%%
%% Minted listings support 
%% Need pygment <http://pygments.org/> <http://pypi.python.org/pypi/Pygments>
\usepackage{listings}
\usepackage{xcolor}

\lstdefinelanguage{JavaScript}{
  keywords={const, let, var, function, return, if, else, for, while, class, new, await, async},
  sensitive=true,
  comment=[l]{//},
  morecomment=[s]{/*}{*/},
  morestring=[b]",
  morestring=[b]'
}

%% auto break lines
\lstset{breaklines=true}

\lstset{
  language=JavaScript,
  basicstyle=\ttfamily\small,
  keywordstyle=\bfseries,
  commentstyle=\itshape,
  stringstyle=\ttfamily,
  numbers=left,
  numberstyle=\tiny,
  frame=single,
  breaklines=true,
  columns=fullflexible
}

% to active note:
%\usepackage{todonotes}
% to unactive note:
\usepackage[disable]{todonotes}

%%
%% end of the preamble, start of the body of the document source.
\begin{document}

%%
%% Rights management information.
%% CC-BY is default license.
\copyrightyear{2022}
\copyrightclause{Copyright for this paper by its authors.
  Use permitted under Creative Commons License Attribution 4.0
  International (CC BY 4.0).}

%%
%% This command is for the conference information
\conference{Woodstock'22: Symposium on the irreproducible science,
  June 07--11, 2022, Woodstock, NY}

%%
%% The "title" command
\title{Unravel-RDF: A JavaScript Framework for Building Web Applications over RDF Knowledge Graphs}

\tnotemark[1]
\tnotetext[1]{You can use this document as the template for preparing your
  publication. We recommend using the latest version of the ceurart style.}

%%
%% The "author" command and its associated commands are used to define
%% the authors and their affiliations.
\author[1,2,3]{Olivier Filangi}[%
orcid=0000-0002-2094-3271,
email=olivier.filangi@inrae.fr,
url=https://forge.inrae.fr/olivier.filangi,
]
\cormark[1]
\fnmark[1]
\address[1]{Institute for Genetics, Environment and Plant Protection (IGEPP), National Research Institute for Agriculture, Food and Environment (INRAE), Université Rennes, Institut Agro, Le Rheu, France}
\address[2]{Metabolic Profiling and Metabolomics platform (P2M2), Institute for Genetics, Environment and Plant Protection (IGEPP), Biopolymers Interactions Assemblies (BIA), Le Rheu, France}
\address[3]{MetaboHUB, National Infrastructure of Metabolomics and Fluxomics, Toulouse, France}

\author[4]{Ilaria Tiddi}[%
orcid=0000-0001-7116-9338,
email=i.tiddi@vu.nl,
url=https://kmitd.github.io/ilaria/,
]
\fnmark[1]
\address[4]{Vrije Universiteit Amsterdam, De Boelelaan 1105, 1081 HV Amsterdam, The Netherlands}

\author[5]{Manfred Jeusfeld}[%
orcid=0000-0002-9421-8566,
email=Manfred.Jeusfeld@acm.org,
url=http://conceptbase.sourceforge.net/mjf/,
]
\fnmark[1]
\address[5]{University of Skövde, Högskolevägen 1, 541 28 Skövde, Sweden}

%% Footnotes
\cortext[1]{Corresponding author.}
\fntext[1]{These authors contributed equally.}

%%
%% The abstract is a short summary of the work to be presented in the
%% article.
\begin{abstract}
  A clear and well-documented \LaTeX{} document is presented as an
  article formatted for publication by CEUR-WS in a conference
  proceedings. Based on the ``ceurart'' document class, this article
  presents and explains many of the common variations, as well as many
  of the formatting elements an author may use in the preparation of
  the documentation of their work.
\end{abstract}

%%
%% Keywords. The author(s) should pick words that accurately describe
%% the work being presented. Separate the keywords with commas.
\begin{keywords}
  LaTeX class \sep
  paper template \sep
  paper formatting \sep
  CEUR-WS
\end{keywords}

%%
%% This command processes the author and affiliation and title
%% information and builds the first part of the formatted document.
\maketitle

%1. Introduction
%2. Architecture
%3. Programming Model / API
%4. Design Patterns
%5. Case Study
%6. Discussion
%7. Related Work
%8. Conclusion

\section{Introduction}

Semantic Web technologies are now widely used to represent and query knowledge in the life sciences\citep{wu_semantic_2014}. 
Many resources are published as RDF knowledge graphs and made available via SPARQL endpoints, enabling interoperability 
between heterogeneous datasets describing molecular entities, biological processes, pathways, and phenotypic characteristics 
across diverse organisms.\citep{wilkinson_fair_2016}. In bioinformatics, these approaches enable the integration and reuse 
of omics resources, but their effective use within analysis platforms and web applications still poses significant technical challenges.

Integrating knowledge graphs into web applications with complex user interfaces remains a significant challenge in modern 
software development. As developers increasingly integrate RDF graph data models into web applications, 
they face a primary obstacle: bridging RDF graph data models and application logic. 
Existing RDF middleware frameworks provide limited support for building applications backed by RDF knowledge graphs, 
which introduces additional implementation complexity.

At the same time, query engines such as Comunica have greatly simplified the execution of queries on decentralised RDF data 
sources from the Web. However, creating interactive applications based on these engines remains a real challenge. 
Developers still have to embed raw SPARQL strings in application code, resulting in a weak integration between query 
construction and software abstractions.

By contrast, software development for relational databases has long benefited from sophisticated programming abstractions, 
such as query builders, domain-specific languages (DSLs), and object-relational mappers (ORMs), which improve readability, 
modularity, and maintainability while preserving access to the underlying query language. Comparable abstractions remain 
limited within the Semantic Web ecosystem, particularly for JavaScript applications targeting RDF data sources.

To address this gap, we introduce \texttt{Unravel-RDF}, a JavaScript framework that combines a domain-specific language (DSL) for 
SPARQL query construction with a composable API enabling modular query building, and supports their execution within event-driven 
user interface applications.

We evaluate and illustrate \texttt{Unravel-RDF} using the FORVM ecosystem, which comprises a reference biomedical knowledge graph 
and an experimental extension targeting plant sciences. We demonstrate a cross-domain exploration workflow in which biomedical 
and plant-related knowledge graphs are interconnected through shared chemical entities, enabling interactive navigation of 
metabolomics relationships across domains. This workflow is implemented as a Web application built with \texttt{Unravel-RDF}, 
which bridges query construction and interface logic by orchestrating SPARQL query composition, reuse, and execution across 
multiple endpoints within an event-driven architecture, while providing an interactive visualization layer for exploring 
heterogeneous RDF data sources.

\section{Architecture}

\begin{figure}[t]
    \centering
    \includegraphics[width=0.50\linewidth]{Unravel-rdf-architecture.png}
    \caption{Unravel-RDF architecture: a layered design connecting RDF data sources, the Unravel query abstraction layer, and UI integration components.}
    \label{fig:architecture}
\end{figure}

\subsection{Overview}

Figure~\ref{fig:architecture} illustrates the architecture of Unravel-RDF, which is organized as a layered system connecting RDF data sources, a compositional query abstraction layer, and event-driven application components.

At its core, Unravel-RDF provides a domain-specific language (DSL) for constructing SPARQL queries programmatically. This DSL replaces raw query strings with structured, composable query fragments that support modular construction and variable scoping. Queries are incrementally assembled and compiled into executable SPARQL expressions.

Query execution is delegated to the Comunica framework, which handles access to distributed RDF data sources. Unravel-RDF does not reimplement query processing; instead, it focuses on bridging high-level query composition with execution through a typed Scala.js facade that wraps the Comunica JavaScript API. This facade ensures type-safe interaction between the Scala-based abstraction layer and the underlying JavaScript runtime.

Finally, query results are exposed through an event-driven execution model. Rather than returning static result sets, Unravel-RDF propagates intermediate and final results as events, enabling tight integration with reactive user interface components and incremental data consumption in web applications.

The remainder of this section details these three layers: query construction, execution integration, and event-driven programming.
These three layers are further decomposed into five main functional components that structure the Unravel-RDF architecture, ranging from low-level query construction primitives to high-level UI integration mechanisms.


\subsection{Query Construction}

Query construction in Unravel-RDF is based on a layered design that replaces raw SPARQL string manipulation with a structured and compositional API. This layer is responsible for transforming modular query fragments into executable SPARQL expressions.

It is composed of four main mechanisms:

\begin{itemize}
\item \textbf{Configuration.} Defines the execution context of an Unravel session, including SPARQL prefixes, named graphs, and RDF data sources.

\item \textbf{Core Query DSL.} Provides composable primitives for constructing SPARQL queries from graph patterns, filters, and variable bindings.

\item \textbf{High-Level Query API.} Introduces reusable abstractions for common RDF querying patterns such as graph traversal and decomposed property retrieval.

\item \textbf{Variable Scoping and Focus.} Introduces a dynamic focus mechanism that structures query construction around a central variable context.
\end{itemize}

Together, these components enable incremental query construction and automatic SPARQL generation while preserving full expressivity.

\subsection{Execution Integration}

The execution layer of Unravel-RDF is responsible for bridging the high-level query construction DSL with the underlying SPARQL execution engine. Rather than implementing query processing internally, Unravel-RDF delegates execution to the Comunica framework, which provides efficient access to distributed RDF data sources.

A key aspect of this layer is the translation of compositional query objects into executable SPARQL queries. Once a query is fully constructed, it is compiled into a standard SPARQL expression and passed to Comunica for evaluation. This design ensures a clear separation between query specification and execution strategy.

To enable seamless interoperability between the Scala-based abstraction layer and the JavaScript runtime of Comunica, Unravel-RDF relies on a Scala.js facade. This facade acts as a typed binding layer that wraps the Comunica API, allowing the Scala implementation to interact with JavaScript modules in a type-safe and structured manner while preserving full compatibility with the underlying execution engine.

During execution, results are not simply returned as static collections. Instead, they are progressively emitted and streamed back to the application layer, enabling incremental consumption and integration with reactive interfaces.

\subsection{Event-driven Programming}

Unravel-RDF adopts an event-driven execution model that tightly integrates query processing with application-level logic. Instead of relying on synchronous result retrieval, query execution produces a stream of events that reflect both intermediate and final results.

This model is particularly suited for web applications built around dynamic and interactive user interfaces, where incremental updates are preferred over blocking query execution. As queries are executed through Comunica, result events are propagated asynchronously to the application layer.

The execution pipeline can be summarized as follows: query construction, execution through Comunica, event emission, and consumption by application components. This allows UI components to react directly to query progression, enabling patterns such as live result rendering, progressive data loading, and interactive exploration of RDF graphs.

The event system also provides hooks for integrating additional application logic, such as result transformation, logging, or enrichment. In combination with the execution integration layer, it forms the basis for a reactive architecture where query execution and interface behavior are decoupled yet tightly synchronized.


\section{Programming Model}

This section illustrates how Unravel-RDF is used in practice to construct and execute SPARQL queries within JavaScript applications. We present concrete examples that demonstrate how the abstractions introduced by the framework translate into executable query workflows.

\subsection*{Focus-based query construction and query incremental building}

Queries are constructed relative to a focus variable, which simplifies
navigation over RDF graphs. In this example, the query is centered on a compound entity .

This enables modular construction of complex federated queries from smaller reusable building blocks.

\begin{lstlisting}[caption={Focus-based query construction example},label={lst:example-1}]

  const configuration = UnravelConfiguration
                        .init()
                        .sparqlEndpoint("https://rdfportal.org/ebi/sparql");

  /* Build a base query with "compound" as the current focus, linked to CHEBI:106243 through the variable ?rel */
  const query = UnravelSession(configMetaNetX)
  .prefix("CHEBI", "http://purl.obolibrary.org/obo/CHEBI_")
  .something(
    "compound",compound => compound.out("?rel","CHEBI:24279") /* SPARQL : ?compound ?rel CHEBI:106243 */
  );
  /* Restore the focus to "subject" and retrieve its RDF type */
  const specializedQuery = query
  .from(
    "compound", compound => compound.out("rdf:type","?type")
  );
\end{lstlisting}


- `something("subject", subject => ...)` creates a **lambda enclosure** around the current focus `subject`.
- Inside the closure, `out("?rel", "CHEBI:106243")` searches for RDF triples where `subject` is connected to `CHEBI:106243` by some relation bound to `?rel`.
- `select("subject", "rel")` turns this traversal into a query projection.
- `limit(3).commit().raw()` executes the resulting query and returns the raw SPARQL JSON response.

\subsection*{Event-driven execution}

Query execution is asynchronous and returns a stream of results.

\begin{verbatim}
### Automatic lazy pagination

No manual OFFSET/LIMIT management:

```js
session
  .commit()
  .raw()
  .selectByPage("study", "label")
  .then(([totalCount, fetchPage]) => {
    console.log(`${totalCount} results`)

    // Fetch page 0 on demand
    fetchPage(0).then(page => renderTable(page))
  })
```

### Query progression and events

Long-running SPARQL queries report progress — useful for updating a loading indicator:

```js
session
  .select("study", "label")
  .commit()
  .progression((percent) => {
    updateProgressBar(percent)
  })
  .requestEvent((event) => {
    console.log("SPARQL event:", event)
  })
  .raw()
  .then(response => render(response))
```

\end{verbatim}

This allows integration with reactive user interfaces such as dashboards or interactive graph explorers.


\section{Case Study: FORVM}

We evaluate Unravel-RDF using the FORVM ecosystem, which provides a cross-domain knowledge graph infrastructure spanning biomedical and plant sciences.

The biomedical component, \textit{FORVM ChemDisease}, is a curated knowledge graph derived from statistical co-occurrence of MeSH-annotated chemical and disease concepts in PubMed literature. In parallel, \textit{FORVM Plants} extends this approach to plant sciences by leveraging ontology-driven annotations (including the Plant Ontology and Trait Ontology) to extract plant-related entities from scientific literature.

These two knowledge graphs are interconnected through shared chemical entities, enabling cross-domain exploration of metabolomics relationships across multiple SPARQL endpoints.

We implement this workflow as a web application built with Unravel-RDF. The system orchestrates compositional query construction, reuse of query fragments, and federated execution across endpoints. The event-driven execution model enables interactive exploration of relationships between biomedical and plant knowledge graphs through incremental result rendering and dynamic interface updates.

\section{Acknowledgments}

Identification of funding sources and other support, and thanks to
individuals and groups that assisted in the research and the
preparation of the work should be included in an acknowledgment
section, which is placed just before the reference section in your
document.

This section has a special environment:
\begin{lstlisting}
\begin{acknowledgments}
  These are different acknowledgments.
\end{acknowledgments}
\end{lstlisting}
so that the information contained therein can be more easily collected
during the article metadata extraction phase, and to ensure
consistency in the spelling of the section heading.

Authors should not prepare this section as a numbered or unnumbered
\verb|\section|; please use the ``\verb|acknowledgments|'' environment.

\section{Appendices}

If your work needs an appendix, add it before the
``\verb|\end{document}|'' command at the conclusion of your source
document.

Start the appendix with the ``\verb|\appendix|'' command:
\begin{lstlisting}
\appendix
\end{lstlisting}
and note that in the appendix, sections are lettered, not
numbered. 

%%
%% The acknowledgments section is defined using the "acknowledgments" environment
%% (and NOT an unnumbered section). This ensures the proper
%% identification of the section in the article metadata, and the
%% consistent spelling of the heading.
\begin{acknowledgments}
  Thanks to the developers of ACM consolidated LaTeX styles
  \url{https://github.com/borisveytsman/acmart} and to the developers
  of Elsevier updated \LaTeX{} templates
  \url{https://www.ctan.org/tex-archive/macros/latex/contrib/els-cas-templates}.  
\end{acknowledgments}

%% The declaration on generative AI comes in effect
%% in Janary 2025. See also
%% https://ceur-ws.org/GenAI/Policy.html
\section*{Declaration on Generative AI}
  {\em Either:}\newline
  The author(s) have not employed any Generative AI tools.
  \newline
  
 \noindent{\em Or (by using the activity taxonomy in ceur-ws.org/genai-tax.html):\newline}
 During the preparation of this work, the author(s) used X-GPT-4 and Gramby in order to: Grammar and spelling check. Further, the author(s) used X-AI-IMG for figures 3 and 4 in order to: Generate images. After using these tool(s)/service(s), the author(s) reviewed and edited the content as needed and take(s) full responsibility for the publication’s content. 

 During the preparation of this work, the authors used Perplexity and ChatGPT in order to: paraphrase
and reword, improve writing style, and grammar and spelling checking. After using these tools, the
authors reviewed and edited the content as needed and take full responsibility for the publication’s
content.

%%
%% Define the bibliography file to be used
\bibliography{sample-ceur}

%%
%% If your work has an appendix, this is the place to put it.
\appendix

\section{Online Resources}


The sources for the ceur-art style are available via
\begin{itemize}
\item \href{https://github.com/yamadharma/ceurart}{GitHub},
% \item \href{https://www.overleaf.com/project/5e76702c4acae70001d3bc87}{Overleaf},
\item
  \href{https://www.overleaf.com/latex/templates/template-for-submissions-to-ceur-workshop-proceedings-ceur-ws-dot-org/pkfscdkgkhcq}{Overleaf
    template}.
\end{itemize}

\end{document}

%%
%% End of file